\documentclass{beamer}
\usetheme{Madrid}

\title{Memory API in UNIX Systems}
\subtitle{Allocating and Managing Memory in C}
\author{}
\date{\today}

\begin{document}

\begin{frame}
    \titlepage
\end{frame}

\begin{frame}[fragile]{Types of Memory in C}
    \begin{itemize}
        \item \textbf{Stack (Automatic) Memory:} Allocations and deallocations are managed implicitly by the compiler. %[cite: 1]
        \item Space is made upon function call and deallocated on return. %[cite: 1]
        \item \textbf{Heap Memory:} All allocations and deallocations are explicitly handled by the programmer, providing a way to store long-lived memory. %[cite: 1]
    \end{itemize}
    \vspace{0.3cm}
    \begin{block}{Code Example: Stack vs. Heap} %[cite: 1]
\begin{verbatim}
// Automatically managed on the stack
int x = 5; 

// Safely allocating on the heap
int *arr = malloc(10 * sizeof(int)); 
\end{verbatim}
    \end{block}
\end{frame}

\begin{frame}[fragile]{Allocating Memory: \texttt{malloc()}}
    \begin{itemize}
        \item The \texttt{malloc(size\_t size)} call takes the number of bytes needed and returns a pointer to the newly-allocated space, or \texttt{NULL} on failure. %[cite: 1]
        \item \textbf{Best Practice:} Use \texttt{sizeof()} as a compile-time operator instead of hardcoding sizes (e.g., \texttt{malloc(40)}). %[cite: 1]
        \item \textbf{String Allocation:} Use \texttt{strlen(s) + 1} to make room for the end-of-string character, as \texttt{sizeof()} will not work correctly for strings. %[cite: 1]
    \end{itemize}
    \vspace{0.3cm}
    \begin{block}{Code Example: Allocation Types} %[cite: 1]
\begin{verbatim}
// Specific allocation for 10 integers
int *arr = malloc(10 * sizeof(int));

// +1 makes room for the null terminator
char *str = malloc(strlen(name) + 1);
\end{verbatim}
    \end{block}
\end{frame}

\begin{frame}[fragile]{Deallocating Memory: \texttt{free()}}
    \begin{itemize}
        \item To free heap memory no longer in use, call \texttt{free()} and pass it the pointer returned by \texttt{malloc()}. %[cite: 1]
        \item The size of the allocated region is tracked internally by the memory-allocation library, so the user does not need to pass it. %[cite: 1]
        \item While the OS reclaims all process memory when a program exits, programmers should maintain the habit of explicitly freeing memory to avoid poor form. %[cite: 1]
    \end{itemize}
    \vspace{0.3cm}
    \begin{block}{Code Example: Proper Deallocation} %[cite: 1]
\begin{verbatim}
int *arr = malloc(10 * sizeof(int));

// Giving the memory back to the system
free(arr); 
\end{verbatim}
    \end{block}
\end{frame}

\begin{frame}[fragile]{Common Memory Errors}
    \begin{itemize}
        \item \textbf{Segmentation Faults:} Passing unallocated pointers to routines like \texttt{strcpy()} causes illegal memory access and crashes. %[cite: 1]
        \item \textbf{Buffer Overflows:} Not allocating enough memory (e.g., forgetting the +1 for a string), which can silently overwrite contiguous blocks of important memory. %[cite: 1]
        \item \textbf{Memory Leaks:} Forgetting to free memory, which eats up RAM and eventually crashes the system. %[cite: 1]
        \item \textbf{Dangling Pointers:} Freeing memory too early or multiple times confuses the allocation library. %[cite: 1]
    \end{itemize}
    \vspace{0.2cm}
    \begin{block}{Code Example: Memory Leak}
\begin{verbatim}
while(1) { 
    malloc(1024); // Classic example of a memory leak
}
\end{verbatim}
    \end{block}
\end{frame}

\begin{frame}[fragile]{Underlying OS Support \& Advanced APIs}
    \begin{itemize}
        \item \texttt{malloc()} and \texttt{free()} are not system calls, but library calls built on top of OS system calls. %[cite: 1]
        \item \textbf{Heap Boundary Control:} System calls \texttt{brk()} and \texttt{sbrk()} extend or shrink the heap of the process in main memory. %[cite: 1]
        \item You should never call these directly; the libraries manage it well. %[cite: 1]
        \item \textbf{Virtual RAM / Swap:} The \texttt{mmap(...)} call can map anonymous memory to swap space to prevent your system from crashing if it runs out of main memory. %[cite: 1]
        \item \textbf{Advanced Library Functions:}
        \begin{itemize}
            \item \texttt{calloc()}: Allocates memory and zeroes it out. %[cite: 1]
            \item \texttt{realloc()}: Resizes an existing memory region. %[cite: 1]
        \end{itemize}
    \end{itemize}
\end{frame}

\end{document}
